% year: תש"פ
% type: בוחן
% semester: א
% points: 20
% source: exams/בחנים/תש״פ, בוחן, 9141.pdf
% number: 4
% categories: derivatives, multiple-choice

\begin{question}
הנגזרת $y'(0)$ של הפונקציה הסתומה $y(x)$ המוגדרת על ידי המשוואה
$\cos\left(\frac{\pi}{2}+x+y\right)-y=0$ בסביבה של הנקודה $(0,0)$ שווה ל-
\begin{enumerate}
  \item $\pi/2$
  \item $-\frac12$
  \item $1+\pi/2$
  \item $\frac12$
  \item $0$
\end{enumerate}
\end{question}

\begin{hint}
גזרו את שני אגפי המשוואה לפי $x$, כאשר $y=y(x)$ (כלל השרשרת), ואז הציבו $x=0,\ y=0$.
\end{hint}

\begin{solution}
\textbf{התשובה הנכונה: (ב)} $-\frac12$.

ראשית, הנקודה $(0,0)$ מקיימת את המשוואה: $\cos\frac{\pi}{2}-0=0$.
נגזור את שני אגפי המשוואה $\cos\left(\frac{\pi}{2}+x+y(x)\right)-y(x)=0$ לפי $x$, תוך שימוש בכלל השרשרת:
\[
-\sin\left(\frac{\pi}{2}+x+y\right)\cdot\left(1+y'\right)-y'=0 .
\]
נציב $x=0$, $y(0)=0$: $\sin\frac{\pi}{2}=1$, ולכן
\[
-(1+y'(0))-y'(0)=0\quad\Longrightarrow\quad -1-2y'(0)=0\quad\Longrightarrow\quad y'(0)=\boxed{-\frac12}.
\]
(בדיקה: לפי הזהות $\cos(\frac{\pi}{2}+\theta)=-\sin\theta$ המשוואה היא $-\sin(x+y)=y$; גזירה נותנת $-\cos(x+y)(1+y')=y'$, ובנקודה $(0,0)$ שוב $y'(0)=-\frac12$.)

\textbf{למה האחרות שגויות:} (ד) מתקבלת משגיאת סימן בנגזרת של $\cos$. (ה) מתקבלת אם שוכחים לגזור את $y$ שבתוך הקוסינוס או את האיבר $-y$. (א) ו-(ג) מתקבלות מבלבול בין ערך הזווית $\frac{\pi}{2}$ לבין הנגזרת; הנגזרת היא ערך יחיד שחישבנו, $-\frac12$.
\end{solution}
